\emph{(i).} The blades are a basis (Theorem \claimref{synthesis.signed-directions-and-the-integer-clifford-algebra}), so every element is the sum of an even and an odd part. Each $e_{j}f_{j}$ is self-adjoint ($(e_{j}f_{j})^{\dagger} = f_{j}^{\dagger}e_{j}^{\dagger} = -f_{j}e_{j} = e_{j}f_{j}$) with $(e_{j}f_{j})^{2} = -Q(e_{j})Q(f_{j}) = 1$, and the $e_{j}f_{j}$ commute, so $\Gamma^{\dagger} = \Gamma$ and $\Gamma^{2} = 1$. A generator anticommutes with the factor $e_{j}f_{j}$ of its own mode and commutes with the others, so it anticommutes with $\Gamma$; hence $\Gamma$ commutes with even elements and anticommutes with odd ones. For $x = x_{0} + x_{1}$ (even and odd parts), $\Gamma x - x\Gamma = 2\Gamma x_{1}$, which vanishes only if $x_{1} = 0$: the even subalgebra is the commutant of $\Gamma$. By the ladder relations (Proposition \claimref{tower.ladder-pairs-satisfy-car-vectors}; here $w_{j}v_{j} = \tfrac12(1 + e_{j}f_{j})$ directly), $e_{j}f_{j} = 2w_{j}v_{j} - 1$ and $w_{j}v_{j}P_{0} = P_{0}$, so $e_{j}f_{j}P_{0} = P_{0}$ and $\Gamma P_{0} = P_{0}$. Moving $\Gamma$ through the $|I|$ vectors of $s_{I} = v_{i_{1}}\cdots v_{i_{m}}P_{0}$ gives $\Gamma s_{I} = (-1)^{|I|}s_{I}$. Hence $s_{I}^{\dagger}\Gamma s_{J} = (-1)^{|J|}\delta_{IJ}P_{0}$ by clause (i) of \claimref{synthesis.additive-weights-on-the-rational-projectors-are-linear-in-every-size-from-four}: $m(\Gamma)$ is the diagonal matrix with entries $(-1)^{|I|}$. Since $m$ is an isomorphism, it maps the commutant of $\Gamma$ onto the commutant of that diagonal matrix, the matrices with $M_{IJ} = 0$ for $|I| + |J|$ odd. There are $2^{\ell-1}$ subsets of each parity, so these are two blocks of size $2^{\ell-1}$, and $m(\tfrac12(1 \pm \Gamma))$ are the block identities. Since $m(x^{\dagger}) = m(x)^{T}$, an even projector has a symmetric idempotent matrix in two blocks, $m(E) = E_{+} \oplus E_{-}$ with $E_{\pm}$ symmetric idempotents, and conversely.

\emph{(ii).} Let $n = 2^{\ell-1} \ge 4$. For a symmetric idempotent $A$ of $\Mat_{n}(\Q)$ put $\mu_{+}(A) = \mu(m^{-1}(A \oplus 0))$ and $\mu_{-}(A) = \mu(m^{-1}(0 \oplus A))$. Both are defined, because $A \oplus 0$ and $0 \oplus A$ are matrices of even projectors by (i), and both are additive weights on the projectors of $\Mat_{n}(\Q)$, because $A \mapsto A \oplus 0$ preserves sums and the relation $AA' = 0$. By clause (iii) of \claimref{synthesis.additive-weights-on-the-rational-projectors-are-linear-in-every-size-from-four} there are unique symmetric $\rho^{\pm}$ with entries in $V$ and $\mu_{\pm}(A) = \sum\rho^{\pm}_{ij}A_{ij}$. For an even projector $E$, $m(E) = E_{+} \oplus E_{-}$ is the sum of the orthogonal projectors $E_{+} \oplus 0$ and $0 \oplus E_{-}$, both even, so $\mu(E) = \mu_{+}(E_{+}) + \mu_{-}(E_{-})$. With $\rho = \rho^{+} \oplus \rho^{-}$ this is the formula of the statement, and $L(x) = \sum_{|I| + |J|\ \mathrm{even}}\rho_{IJ}m(x)_{IJ}$ is linear on the even subalgebra. A linear $L'$ on the even subalgebra with $L'(x^{\dagger}) = L'(x)$ is determined by its values on self-adjoint even elements, whose matrices are the symmetric matrices in two blocks; these are spanned by the matrices of even projectors (in each block, the projectors onto the coordinate rays and onto the rays of sums of two coordinate vectors). So $L' = L$ if $L'$ restricts to $\mu$.

\emph{(iii).} $\mu \ge 0$ on the projectors of the rays of each block gives $\sum\rho^{\pm}_{ij}r_{i}r_{j} \ge 0$ for every rational $r$. For any $x$ of the algebra let $M = m(x)$. $m(\tfrac12(x + \Gamma x\Gamma))$ is the block-diagonal part of $M$, because conjugation by the diagonal matrix $m(\Gamma)$ changes the sign of the entries with $|I| + |J|$ odd; so $\omega$ is linear, vanishes on odd elements, equals $L$ on even ones, and $\omega(x^{\dagger}) = \omega(x)$ by the symmetry of $\rho$. The block-diagonal part of $MM^{T}$ has the blocks $\sum_{k}c^{\pm}_{k}(c^{\pm}_{k})^{T}$, where $c^{+}_{k}$ and $c^{-}_{k}$ are the two halves of the $k$-th column of $M$, so $\omega(xx^{\dagger}) = \sum_{k}\bigl(\sum\rho^{+}_{ij}(c^{+}_{k})_{i}(c^{+}_{k})_{j} + \sum\rho^{-}_{ij}(c^{-}_{k})_{i}(c^{-}_{k})_{j}\bigr) \ge 0$; for even $x$ this is $L(xx^{\dagger}) \ge 0$. $\omega(1) = \mu(1) = 1$. A state that vanishes on odd elements is determined by its restriction to the even subalgebra, which satisfies $L'(x^{\dagger}) = L'(x)$ and restricts to $\mu$, so it is $\omega$ by the uniqueness in (ii).

\emph{(iv).} For $\ell = 1$ the even subalgebra is spanned by the orthogonal idempotents $\tfrac12(1 \pm \Gamma)$, its projectors are the four named, and an additive $\mu$ has $\mu(0) = 0$ and $\mu(1) = \mu(\tfrac12(1 + \Gamma)) + \mu(\tfrac12(1 - \Gamma))$, so it is the restriction of the linear map with those two values. For $\ell = 2$ the blocks have size two. In $\Q^{2}$ the only ray orthogonal to that of $(a, b)$ is that of $(-b, a)$, and a projector of a block is $0$, a ray or the block identity, so an additive weight on the even projectors is the same thing as a pair of functions on the rays of the two blocks, each with $f(a, b) + f(-b, a)$ constant. For a primitive integer vector $(a, b)$ the parities of $(a, b)$ and of $(-b, a)$ are $(1,0)$ and $(0,1)$, or both $(1,1)$, and the sign of the vector does not change them, so the function of the statement is well defined on rays and has $f(a, b) + f(-b, a) = 1$. If it were $\sum\rho_{ij}r_{i}r_{j}/(r \cdot r)$, the values $1$, $0$, $\tfrac12$ at $(1,0)$, $(0,1)$, $(1,1)$ would give $\rho_{11} = 1$, $\rho_{22} = 0$, $\rho_{12} = 0$, hence the value $\tfrac15$ at $(1, 2)$, where the function is $1$.

\emph{The example in (iii).} Let $y = s_{I} + s_{J}$ with $|I|$ even and $|J|$ odd. Then $y^{\dagger}y = 2P_{0}$, and $\omega'(x) = \tfrac12\sum_{A, B \in \{I, J\}}m(x)_{AB}$ is the functional $x \mapsto \tfrac12 a^{T}m(x)a$ with $a$ the sum of the two coordinate vectors: it is linear, $\omega'(1) = 1$, $\omega'(xx^{\dagger}) = \tfrac12|m(x)^{T}a|^{2} \ge 0$, $\omega'(x^{\dagger}) = \omega'(x)$, on even $x$ it is $\tfrac12(m(x)_{II} + m(x)_{JJ})$ by (i), and on the odd element $s_{I}s_{J}^{\dagger} + s_{J}s_{I}^{\dagger}$ it is $1$.

\emph{(v).} Let $\mu$ be such a weight. By (ii) with $V = \Q$ it is linear, $\mu = \mu_{+} + \mu_{-}$. Since $\mu(1) = \mu(\tfrac12(1 + \Gamma)) + \mu(\tfrac12(1 - \Gamma)) = 1$ with both terms in $\{0, 1\}$, one block has value $1$; on it the function $q(r) = \sum\rho_{ij}r_{i}r_{j}/(r \cdot r)$ on rays takes only the values $0$ and $1$ and sums to $1$ on every frame of the block. Take a frame; exactly one of its rays, $r$, has $q(r) = 1$, and another, $s$, has $q(s) = 0$; $r \cdot s = 0$. For rational $t$ put $N(t) = \sum\rho_{ij}(r + ts)_{i}(r + ts)_{j}$ and $D(t) = (r + ts) \cdot (r + ts) = r \cdot r + t^{2}\,s \cdot s > 0$, polynomials of degree at most two. For every rational $t$, $N(t) = 0$ or $N(t) = D(t)$, so the polynomial $N(N - D)$ has infinitely many zeros and vanishes identically, and $N = 0$ or $N = D$ as polynomials. $N(0) = r \cdot r \ne 0$ excludes the first; the coefficient of $t^{2}$ is $\sum\rho_{ij}s_{i}s_{j} = 0$ in $N$ and $s \cdot s > 0$ in $D$, which excludes the second. Three values of $t$ already suffice: with $N(t) = a + 2\beta t$, $a = r \cdot r$, and $D(t) = a + \sigma t^{2}$, $\sigma = s \cdot s$, the conditions at $t = 1$ and $t = -1$ force $a = \sigma$ and $\beta = \pm\tfrac12a$, and then $N(2) = 3a$ or $-a$, neither $0$ nor $D(2) = 5a$; the check runs this on a grid.
